% !PRINTED 245
\weberpage{245}
%<<CONT
cannot, however, be $\equiv 0\ (\mathrm{mod}\ p)$; for otherwise a congruence could be derived from (6):

\[
P \equiv T P_{\nu-1}\ (\mathrm{mod}\ p),
\]

in which $T$ is a non-constant function of $t$, and $P$ would not be irreducible with respect to the modulus $p$.

Now if $Q$ is any integral function of $t$ that satisfies the condition

\[
Q P_1 \equiv 0\ (\mathrm{mod}\ P, p)
\]

then, on multiplying (6) by $Q$, it follows that:

\[
Q P_2 \equiv 0,\ Q P_3 \equiv 0,\ \ldots,\ Q P_\nu \equiv 0\ (\mathrm{mod}\ P, p),
\]

hence $Q \equiv 0$. And if we now take for $P_1$, $Q$ the residues of the functions $F(t)$, $F_1(t)$, then Theorem 1. is hereby proved.

The residues of all functions $F(t)$ with respect to the modulus $P$ are of the form

\weq{7}{X = X(t) = x_0 + x_1 t + \cdots + x_{n-1} t^{n-1},}

and if one wishes to have only those among them that are incongruent with respect to the modulus $p$, it suffices to let the coefficients $x_0, x_1, \ldots x_{n-1}$ run through the sequence of numbers $0, 1, 2, \ldots, p-1$.

\gesperrt{There are therefore $p^n$ and no more functions incongruent with respect to the moduli $P$, $p$.}

\S.\,64.

\gesperrt{Congruence fields.}

It is only another mode of expression for the considerations carried out in the preceding section if, following \gesperrt{Galois}, one introduces an imaginary number $\varepsilon$ which satisfies the congruence

\weq{1}{P(\varepsilon) \equiv 0\ (\mathrm{mod}\ p)}

There is no integral rational number of this kind as soon as $n > 1$, because otherwise $P(t) - P(\varepsilon) = (t-\varepsilon)Q$ would be divisible by $t-\varepsilon$, hence $P(t) \equiv (t-\varepsilon)Q$, and $P(t)$ would not be irreducible. Nevertheless one can, just as in ordinary number theory one introduces the imaginary unit $i = \sqrt{-1}$, make use of such a symbol $\varepsilon$ in calculation. One calculates with it, so far as addition, subtraction and multiplication are concerned, according to the rules of literal calculation, and in doing so may use the equation $P(\varepsilon) = 0$ at pleasure. One can even take $\varepsilon$ to be simply a root of the equation $P(x) = 0$
%CONT>>
